\chapter{Methodology}

\section{Baseline Black-Scholes Pricing}

The analytical Black-Scholes formula is the reference pricing method.
It is used to generate supervised learning labels and to evaluate prediction errors.
Because the formula is available in closed form, it provides an exact ground truth for the controlled experiment.

\section{Monte Carlo Benchmark}

Monte Carlo simulation is used as a numerical benchmark.
For each option, terminal prices are sampled from the exact Black-Scholes terminal distribution and discounted payoffs are averaged.
The benchmark is evaluated against the same analytical Black-Scholes reference as the neural network.

The final Monte Carlo benchmark uses 50,000 simulated paths on a deterministic subset of 512 test options.
The subset restriction keeps the benchmark computationally bounded while still providing a meaningful comparison with a classical simulation-based pricing method.

\section{Neural Network Model}

The neural network is a feed-forward MLP for regression.
The selected final configuration uses:

\begin{itemize}
    \item input features $(S_0, K, T, r, \sigma, S_0/K)$;
    \item hidden layers of widths $64$, $64$, and $32$;
    \item SiLU activation functions in the hidden layers;
    \item one scalar output representing the predicted call price.
\end{itemize}

The optimized loss is the mean squared error:

\begin{equation}
    \mathcal{L}(\theta) =
    \frac{1}{n}
    \sum_{i=1}^{n}
    \left(f_\theta(x_i) - y_i\right)^2.
\end{equation}

\section{Training Procedure}

Inputs are standardized using \texttt{StandardScaler}.
The target is also standardized during training to improve numerical stability and convergence.
Predictions are transformed back to original price units before evaluation.

The network is trained with Adam, mini-batches, weight decay, and early stopping based on validation loss.
The final experiment uses up to 200 epochs and a batch size of 1024.
After prediction, prices are clipped at zero to avoid financially invalid negative option prices.

\section{Evaluation Metrics}

The main metrics are:

\begin{itemize}
    \item Mean Absolute Error (MAE);
    \item Root Mean Squared Error (RMSE);
    \item coefficient of determination $R^2$;
    \item MAPE computed only on prices greater than 1.
\end{itemize}

The MAPE restriction avoids unstable percentage errors for options whose theoretical price is close to zero.

\section{Runtime Benchmark}

The runtime benchmark separates training time from pricing time.
This distinction is important because a neural network has an upfront training cost, while its surrogate value is measured during repeated inference.

The benchmark compares:

\begin{itemize}
    \item vectorized analytical Black-Scholes pricing;
    \item neural network inference after training;
    \item Monte Carlo pricing.
\end{itemize}

\section{SVR Baseline}

The project includes a reduced-scale Support Vector Regression baseline with an RBF kernel.
The SVR benchmark uses 5,000 synthetic samples per run and is repeated over three random seeds.
The same engineered feature set used by the final neural network, including moneyness, is used for SVR.

Both input features and target prices are standardized before fitting.
Predictions are transformed back to original price units before computing metrics.
SVR is kept on reduced datasets because kernel methods scale less favorably than neural networks.

\section{Robustness Experiments with Noisy Targets}

The robustness experiments perturb only the training labels, not the input variables.
Models are trained on noisy Black-Scholes prices and evaluated against clean analytical Black-Scholes prices.

The neural network noisy-target experiment uses 50,000 synthetic contracts per noise level.
The same noisy-target mechanism is also applied to the reduced-scale SVR baseline, using 5,000 contracts per run and three random seeds.

These experiments are controlled label-noise tests.
They should not be interpreted as real market-price modeling, because real exchange-traded options include additional effects such as bid-ask spreads, liquidity, dividends, and implied-volatility structure.
